\section{Conjunction means two obligations}\label{ch15:sec:conjunction}
To prove a conjunction $P\land Q$ you must prove two things, $P$ and $Q$ (Section~\ref{ch02:sec:connectives}). To use a hypothesis $P\land Q$, you may take either half of it. The next theorem does both. From $P\land Q$ it proves $Q\land P$, by taking the two halves apart and putting them back together in the opposite order. On paper: assume that $P$ and $Q$ both hold; then $Q$ holds and $P$ holds, so $Q\land P$.
\par\noindent\begin{minipage}{\linewidth}
\begin{lstlisting}
theorem and_swap (P Q : Prop) : And P Q -> And Q P := by
  intro h
  exact And.intro h.right h.left
\end{lstlisting}
\end{minipage}\par

Lean writes the conjunction $P\land Q$ as \code{And P Q}, and the infoview displays it as $P\land Q$. After \code{intro h}, the fact \code{h : And P Q} is available and the goal is \code{And Q P}. A proof of a conjunction is a pair of proofs, one for each half, and each half can be taken out of the pair: \code{h.left} is a proof of $P$, the first component, and \code{h.right} is a proof of $Q$, the second. In the other direction, \code{And.intro} builds a pair: if \code{a} is a proof of $R$ and \code{b} is a proof of $S$, then \code{And.intro a b} is a proof of \code{And R S}. The order of the two arguments must match the conjunction being built. The goal \code{And Q P} needs a proof of $Q$ first and a proof of $P$ second, so the last line is \code{And.intro h.right h.left}.

The proof above supplies both halves at once. A second version makes the two obligations visible as two separate goals:
\par\noindent\begin{minipage}{\linewidth}
\begin{lstlisting}
theorem and_swap_two_goals (P Q : Prop) :
    And P Q -> And Q P := by
  intro h
  constructor
  exact h.right
  exact h.left
\end{lstlisting}
\end{minipage}\par
The tactic \code{constructor} replaces the goal \code{And Q P} by two goals, $Q$ and $P$, which the infoview labels \code{left} and \code{right}. Each \code{exact} closes the first goal that is still open, so the first one proves $Q$ and the second proves $P$. If the last line is left out, Lean reports the error ``unsolved goals'' and displays the goal $P$ that is still open.

This is the formal counterpart of a paper proof that must establish two set containments (Section~\ref{ch03:sec:sets}), or both directions of an ``if and only if.'' After a long first half it is easy to feel that the proof is finished. Lean keeps track of every open goal, so it does not let that mistake pass, and on paper we must keep the same list ourselves.
